% ============================================================
% Section 60: 单质晶体空位缺陷的振动熵修正
% ============================================================
\newpage
\section{固体理论：单质晶体空位缺陷的振动熵修正}
\label{sec:vacancy-vibrational-entropy}

\begin{modelbox}[题目：考虑晶格振动影响的空位平衡浓度]
在由 $N$ 个原子组成的单原子简单立方晶体中，考虑处于热平衡的空位形成的点缺陷。
\begin{enumerate}
    \item 若晶体的温度为 $T$，产生一个空位所需要的能量为 $u$，写出热平衡时这种缺陷的数目 $n$（$n\ll N$）；
    \item 现在考虑晶格振动（即声子）对结果的影响，即把每个原子看作是独立的振子，对于一个给定的原子，视其周围有无一个空位来选择不同的频率 $\omega'$ 和 $\omega$，此时，缺陷的数目 $n$ 应做何修正？
    \item 两个频率中哪一个大？
\end{enumerate}
\end{modelbox}

\begin{tipbox}[考点快速分析]
\begin{itemize}
    \item \textbf{单质肖特基缺陷}：仅一类空位，$W=\binom{N}{n}$，经典结果 $n=N e^{-u/k_{\!B}T}$（指数无 $1/2$）
    \item \textbf{振动熵修正}：空位改变近邻原子的局域环境，导致振动频率软化/硬化，引入附加的振动自由能
    \item \textbf{爱因斯坦模型}：每个原子视为独立谐振子，高温极限 $k_{\!B}T\gg\hbar\omega$ 下自由能 $f=k_{\!B}T\ln(\hbar\omega/k_{\!B}T)$
    \item \textbf{结构因子}：简单立方中一个空位有 6 个最近邻，每个原子 3 个自由度，故振动模式变化数为 $6\times 3=18$
\end{itemize}
\end{tipbox}

\begin{derivationbox}[标准解题过程]
\textbf{Step 1: 仅考虑位形熵的空位平衡浓度}

从 $N$ 个格点中移走 $n$ 个原子形成 $n$ 个空位，可能的方式数为
\begin{equation}
W = \binom{N}{n} = \frac{N!}{(N-n)!\,n!}.
\end{equation}
组态熵为 $\Delta S = k_{\!B}\ln W$。产生 $n$ 个空位使内能增加 $\Delta U = nu$。自由能变化为
\begin{equation}
\Delta F = \Delta U - T\Delta S = nu - k_{\!B}T\ln\!\left[\frac{N!}{(N-n)!\,n!}\right].
\end{equation}
利用斯特林近似 $\ln x! \approx x\ln x - x$，并假设 $n\ll N$，对 $n$ 求导令其为零：
\begin{equation}
\frac{\partial\Delta F}{\partial n} = u - k_{\!B}T\ln\!\left(\frac{N}{n}\right) = 0.
\end{equation}
解得经典结果
\begin{equation}
\boxed{n = N e^{-u/k_{\!B}T}.}
\label{eq:vacancy-classical}
\end{equation}

\textbf{Step 2: 考虑振动频率变化的修正}

\textit{振动自由能（爱因斯坦模型）。}在高温近似（$\hbar\omega\ll k_{\!B}T$）下，一个频率为 $\omega$ 的谐振子对自由能的贡献为
\begin{equation}
f(\omega) = k_{\!B}T\ln\!\left(\frac{\hbar\omega}{k_{\!B}T}\right).
\end{equation}
（推导：配分函数 $Z=\sum_n e^{-(n+1/2)\hbar\omega/k_{\!B}T}\approx k_{\!B}T/\hbar\omega$，则 $F=-k_{\!B}T\ln Z=k_{\!B}T\ln(\hbar\omega/k_{\!B}T)$。）

\textit{无缺陷时的振动自由能。}完整晶体有 $N$ 个原子，每个原子 3 个振动自由度，总模式数 $3N$。设所有原子振动频率均为 $\omega$，则
\begin{equation}
F_{\mathrm{vib}}^{(0)} = 3N k_{\!B}T\ln\!\left(\frac{\hbar\omega}{k_{\!B}T}\right).
\end{equation}

\textit{有空位时的振动自由能。}在简单立方晶体中，每个空位有 6 个最近邻原子。这些原子由于失去了一个近邻键，所受恢复力减弱，振动频率变为 $\omega'$。
\begin{itemize}
    \item 空位最近邻原子数：$6n$ 个，每个 3 个自由度，共 $18n$ 个模式，频率为 $\omega'$
    \item 其余原子数：$N-6n$ 个，每个 3 个自由度，共 $3(N-6n)$ 个模式，频率仍为 $\omega$
\end{itemize}
有空位时的总振动自由能为
\begin{equation}
F_{\mathrm{vib}} = 18n k_{\!B}T\ln\!\left(\frac{\hbar\omega'}{k_{\!B}T}\right) + 3(N-6n)k_{\!B}T\ln\!\left(\frac{\hbar\omega}{k_{\!B}T}\right).
\end{equation}

\textit{振动自由能的变化。}
\begin{align}
\Delta F_{\mathrm{vib}}
&= F_{\mathrm{vib}} - F_{\mathrm{vib}}^{(0)} \notag\\
&= 18n k_{\!B}T\ln\!\left(\frac{\hbar\omega'}{k_{\!B}T}\right) + 3(N-6n)k_{\!B}T\ln\!\left(\frac{\hbar\omega}{k_{\!B}T}\right) - 3N k_{\!B}T\ln\!\left(\frac{\hbar\omega}{k_{\!B}T}\right) \notag\\
&= 18n k_{\!B}T\left[\ln\!\left(\frac{\hbar\omega'}{k_{\!B}T}\right) - \ln\!\left(\frac{\hbar\omega}{k_{\!B}T}\right)\right] \notag\\
&= 18n k_{\!B}T\ln\!\left(\frac{\omega'}{\omega}\right).
\label{eq:vib-free-energy-change}
\end{align}

\textit{总自由能变化与平衡浓度。}总自由能变化为
\begin{equation}
\Delta F = nu - k_{\!B}T\ln\!\left[\frac{N!}{(N-n)!\,n!}\right] + 18n k_{\!B}T\ln\!\left(\frac{\omega'}{\omega}\right).
\end{equation}
对 $n$ 求导并令其为零（利用斯特林公式）：
\begin{equation}
u - k_{\!B}T\ln\!\left(\frac{N}{n}\right) + 18k_{\!B}T\ln\!\left(\frac{\omega'}{\omega}\right) = 0.
\end{equation}
解得修正后的平衡浓度
\begin{equation}
\boxed{n = N\left(\frac{\omega}{\omega'}\right)^{\!18} e^{-u/k_{\!B}T}.}
\label{eq:vacancy-vib-corrected}
\end{equation}

\textit{修正因子的物理意义。}因子 $(\omega/\omega')^{18}$ 来源于空位周围 $18$ 个振动模式的频率变化（$18=6$ 个近邻 $\times 3$ 个自由度）。振动频率的改变相当于引入了一个附加的振动熵，从而改变了缺陷的\textbf{有效形成自由能}：
\begin{equation}
u_{\mathrm{eff}} = u - 18k_{\!B}T\ln\!\left(\frac{\omega}{\omega'}\right).
\end{equation}

\textbf{Step 3: 频率大小的比较}

当原子周围存在空位时，该原子一侧失去了近邻原子的排斥和吸引作用，等效恢复力减弱，因此其振动频率降低。即
\begin{equation}
\boxed{\omega' < \omega.}
\end{equation}
因此 $\omega/\omega' > 1$，振动熵的贡献使空位的表观浓度比仅考虑形成能 $u$ 时\textbf{更大}。这一效应在高温下尤其显著，因为振动熵项与 $T$ 成正比，而位形熵的``有效能'' $k_{\!B}T\ln(N/n)$ 也随 $T$ 增大。
\end{derivationbox}

\begin{formulabox}[答案汇总]
\begin{center}
\begin{tabular}{ll}
\toprule
\textbf{问题} & \textbf{答案} \\
\midrule
(1) 空位数目（无振动修正） & $n = N e^{-u/k_{\!B}T}$ \\
(2) 考虑振动频率修正 & $\displaystyle n = N\left(\frac{\omega}{\omega'}\right)^{\!18} e^{-u/k_{\!B}T}$ \\
(3) 频率大小关系 & $\omega' < \omega$（空位近邻原子振动频率降低） \\
\bottomrule
\end{tabular}
\end{center}
\end{formulabox}

\begin{insightbox}[物理图像：振动熵如何``帮助''缺陷形成？]
空位不仅增加\textbf{位形熵}（原子可以挪位置），还增加\textbf{振动熵}（近邻原子``松绑''后振动更自由）。

从自由能角度看：
\begin{itemize}
    \item 形成能 $u$：$\Delta U = nu$（$>0$，阻碍缺陷形成）
    \item 位形熵：$-T\Delta S_{\mathrm{config}} = -k_{\!B}T\ln W$（$<0$，促进缺陷形成）
    \item 振动熵：$\Delta F_{\mathrm{vib}} = 18nk_{\!B}T\ln(\omega'/\omega)$（$<0$，因为 $\omega'<\omega$，进一步促进缺陷形成）
\end{itemize}
\textbf{有效形成能} $u_{\mathrm{eff}} = u - 18k_{\!B}T\ln(\omega/\omega') < u$。振动软化相当于给缺陷形成提供了``熵补贴''，使高温下缺陷浓度远高于仅用 $e^{-u/k_{\!B}T}$ 估算的值。
\end{insightbox}

\begin{thinkbox}[考场速记：振动熵修正的通用模板]
遇到``考虑晶格振动对缺陷浓度的影响''类问题：
\begin{enumerate}
    \item \textbf{确定位形熵}：写出 $W(n)$，经典结果 $n_0\sim e^{-u/k_{\!B}T}$
    \item \textbf{数振动模式变化}：
    \begin{itemize}
        \item 确定缺陷的近邻原子数 $z$（简单立方 $z=6$，fcc $z=12$，bcc $z=8$）
        \item 每个原子自由度 $d=3$
        \item 总变化模式数 $=z\times d\times n$
    \end{itemize}
    \item \textbf{写振动自由能变化}：$\Delta F_{\mathrm{vib}} = zdn\,k_{\!B}T\ln(\omega'/\omega)$
    \item \textbf{合并求极值}：有效形成能 $u_{\mathrm{eff}} = u - zdk_{\!B}T\ln(\omega/\omega')$
    \item \textbf{判断频率变化}：缺键$\Rightarrow$恢复力减弱$\Rightarrow$\textbf{频率降低}（$\omega'<\omega$）
\end{enumerate}
\textit{陷阱提醒}：振动熵修正只在高温极限（$k_{\!B}T\gg\hbar\omega$）下成立。低温时需保留零点能 $+\frac{1}{2}\hbar\omega$，修正形式不同。
\end{thinkbox}
